% Part: computability
% Chapter: computability-theory
% Section: applications-fixed-point

\documentclass[../../../include/open-logic-section]{subfiles}

\begin{document}

\olfileid[fr]{cmp}{thy}{apf}
\olsection{Applications du théorème du point fixe}

Le théorème du point fixe permet, pour l'essentiel,
de définir des fonctions calculables partielles
à partir de leurs propres indices. Par exemple,
on peut trouver un indice~$e$ tel que, pour
tout~$y$,
\[
\cfind{e}(y) = e + y.
\]
La preuve du théorème permet aussi de construire,
en Java ou en C++, un programme qui affiche
son propre texte.

Rappelons que $W_e$ désigne le domaine de
$\cfind{e}$ et que la suite $W_0$, $W_1$,
$W_2$,~\dots énumère les ensembles effectivement
énumérables. Certains de ces ensembles sont
calculables. Existe-t-il un algorithme qui,
à partir d'un indice~$x$, renvoie un indice
de la fonction caractéristique de~$W_x$
lorsque cet ensemble est calculable ?
La réponse est négative.

\begin{thm}
Il n'existe aucune fonction calculable partielle~$f$
ayant la propriété suivante : chaque fois que
$W_e$ est calculable, $f(e)$ est définie et
$\cfind{f(e)}$ est la fonction caractéristique
de~$W_e$.
\end{thm}

\begin{proof}
Supposons, pour obtenir une contradiction, qu'une
fonction calculable partielle~$f$ possède cette
propriété. Le théorème du point fixe fournit un
indice~$e$ tel que
\[
\cfind{e}(y) \simeq
\begin{cases}
0 & \text{si $y=0$ et $\cfind{f(e)}(0) \fdefined = 0$} \\
\text{indéfinie} & \text{sinon.}
\end{cases}
\]
En effet, on l'applique à la fonction
\[
g(x,y) \simeq
\begin{cases}
0 & \text{si $y=0$ et $\cfind{f(x)}(0) \fdefined = 0$} \\
\text{indéfinie} & \text{sinon.}
\end{cases}
\]
La fonction~$g$ est calculable partielle :
sur l'entrée $(x,y)$, on vérifie d'abord si
$y=0$. Si oui, on calcule $f(x)$, puis on
utilise la fonction universelle pour calculer
$\cfind{f(x)}(0)$. Si ce dernier calcul
s'arrête avec la valeur~$0$, on renvoie~$0$.
Dans tous les autres cas, y compris lorsque
$f(x)$ n'est pas définie, la procédure ne
s'arrête pas.

Si $f(e)$ est définie et si
$\cfind{f(e)}(0)\fdefined=0$, alors
$\cfind{e}$ n'est définie qu'en~$0$ et
$W_e=\{0\}$. Sinon, $W_e=\emptyset$.
Dans les deux cas, $W_e$ est calculable.
L'hypothèse sur~$f$ impose donc que $f(e)$
soit définie et que $\cfind{f(e)}$ soit la
fonction caractéristique de~$W_e$. Or, dans
le premier cas, elle répond~$0$ à l'entrée~$0$
alors que $0\in W_e$; dans le second, elle ne
répond pas~$0$ à cette entrée alors que
$0\notin W_e$. C'est contradictoire.\footnote{Correction éditoriale : la source commence la preuve avec une fonction $f$ calculable totale, alors que l'énoncé exclut aussi les fonctions calculables partielles. Nous conservons cette portée et incluons le cas où $f(e)$ est indéfinie.}
\end{proof}

\end{document}
